\chapter{Introduction}
\label{ch:introduction}

\section{Motivation}
\label{sec:intro-motivation}

Large language models (LLMs) are increasingly used for multi-sentence generation
tasks such as news summarization. In this setting, the perceived quality of an output is often
dominated not by its overall fluency but by a small number of high-impact factual
errors: a wrong entity, a flipped relation, or a fabricated claim can outweigh
otherwise well-written text. Throughout this thesis, \emph{reliability} refers
specifically to factual consistency with the source document, not to stylistic
polish or fluency.

Preference-based fine-tuning, most prominently reinforcement learning from human
feedback (RLHF), updates a model using comparative judgments between candidate
outputs rather than token-level supervision \citep{christiano2017deep,
stiennon2020learning, ouyang2022training}. Collecting high-quality human preferences
at scale is expensive and slow, so reinforcement learning from AI feedback (RLAIF)
replaces human annotators with an automatic judge \citep{bai2022constitutional,
lee2023rlaif}. This trades annotation cost for a new source of label noise and bias.

A further question arises for long, multi-sentence outputs such as news summaries,
one that is largely independent of whether the feedback comes from a human or an AI
judge: how granular should the preference signal be? A single holistic judgment
(``summary A is better than summary B'') collapses all of the evidence in a
multi-sentence output into one binary decision. If a summary is mostly accurate but
contains one fabricated sentence, a holistic judge may still prefer it over an
alternative that is uniformly mediocre but contains no outright fabrication. The
holistic label cannot express that distinction. Sentence-level judging can in
principle localize exactly which claims are, and are not, supported by the source,
and that localized evidence can then be aggregated into a summary-level decision.
This thesis asks whether that extra granularity actually produces a better training
signal, and under what conditions.

\section{Problem Statement}
\label{sec:intro-problem}

When constructing AI-generated preference data for training a factuality-oriented
reward model, it is unclear whether coarse, holistic preferences are sufficient, or
whether they systematically underuse the supervision signal available in long,
multi-sentence outputs. The problem addressed here, under controlled conditions, is
whether sentence-level factuality judgments, aggregated into a summary-level
preference through a mechanism referred to as \emph{Granular Credit Assignment}
(GCA), yield a more learnable preference signal for reward-model training than
purely holistic judgments of the same candidate summaries.

\begin{definitionbox}[Learnability, as used throughout this thesis]
The pairwise validation accuracy of a Bradley--Terry reward model trained on
the resulting preference pairs (Chapter~\ref{ch:background},
Section~\ref{sec:bg-bradley-terry}): given a held-out preference pair, how
often does the trained reward model rank the summary the judge preferred
above the alternative?
\end{definitionbox}

\section{Thesis Objective}
\label{sec:intro-objective}

The objective of this thesis is to determine, through a controlled empirical
comparison, whether Granular Credit Assignment produces a more learnable
factuality-preference signal than holistic scoring. The two conditions differ only
in how factuality scores are obtained and converted into a pairwise label:

\begin{quote}
\textbf{Holistic:} article $+$ complete candidate summary $\rightarrow$ factuality
score $\rightarrow$ pairwise preference.

\textbf{GCA:} article $+$ individual sentences $\rightarrow$ sentence factuality
scores $\rightarrow$ aggregation $\rightarrow$ pairwise preference.
\end{quote}

Both preference sets are then used to train the same Bradley--Terry reward-model
architecture under the same procedure. Holding the source articles, the candidate
summaries and their generation settings, the evaluator family, the reward-model
architecture, and the training procedure fixed allows feedback granularity to be
studied as the primary experimental variable, rather than as one difference among
several.

The thesis does not propose a new reward-model architecture, a new factuality
metric, or a new training objective. Its subject is the construction of the
preference data itself.

\section{Research Questions}
\label{sec:intro-rqs}

Four research questions follow from this objective. They are stated here in
summary form; Chapter~\ref{ch:methodology},
Section~\ref{sec:methodology-rqs} gives them together with the working hypotheses
and the scope conditions under which each can be answered, and
Chapter~\ref{ch:conclusion}, Section~\ref{sec:conclusion-rqs} answers them
against the evidence collected.

\begin{itemize}
  \item \textbf{RQ1.} Does sentence-level factuality feedback aggregated via GCA
    produce a more learnable pairwise preference signal than holistic
    summary-level feedback, as measured by reward-model pairwise validation
    accuracy?

  \item \textbf{RQ2.} How do evaluator configuration and sentence-score
    aggregation strategy affect the relative learnability of holistic and
    GCA-derived preferences?

  \item \textbf{RQ3.} How stable are the observed differences across repeated
    reward-model training runs and dataset sizes?

  \item \textbf{RQ4.} Does the learnability advantage measured in RQ1 also hold
    when the trained reward models are evaluated against an independent,
    held-out ground truth, rather than against their own training-style
    labels, and does that hold as the training-set size grows?
\end{itemize}

RQ1 is the confirmatory question and is addressed by a repeated-run comparison.
RQ2 is addressed by an exploratory sweep, and therefore supports statements about
association between configuration and outcome rather than about causal mechanism.
RQ3 motivates both the repeated runs at a single dataset size and the
multi-scale check at $n=5{,}000$ and $n=10{,}000$. RQ4 was added after that
multi-scale check, to test the learnability finding against data the trained
models never saw and never influenced their selection
(Chapter~\ref{ch:methodology}, Section~\ref{sec:methodology-ground-truth}).

\section{Contributions}
\label{sec:intro-contributions}

This work makes the following contributions:

\begin{enumerate}
  \item A controlled framework for comparing holistic and sentence-level,
    GCA-based construction of factuality preferences for abstractive
    summarization, in which the candidate pool, evaluator, reward-model
    architecture, and training procedure are held fixed
    (Chapter~\ref{ch:methodology}).

  \item An empirical study of how evaluator mode and aggregation strategy
    interact with feedback granularity. Neither the aggregation formula nor the
    evaluator configuration is neutral: a penalty-weighted aggregation performs
    worse than a simple mean, and in an exploratory sweep the evaluator mode is
    associated with whether an advantage for GCA is observed at all
    (Chapter~\ref{ch:results}, Sections~\ref{sec:results-formula-opt}
    and~\ref{sec:results-mode-sweep}).

  \item A repeated-run, multi-scale evaluation establishing a statistically
    significant GCA advantage at $n=1{,}000$
    (20 independent runs, mean $+3.95$ percentage points, run-level
    Wilcoxon $p<0.001$; found during final review to have been trained
    under AlignScore's \texttt{nli\_sp} mode rather than the \texttt{nli}
    mode intended, Chapter~\ref{ch:implementation},
    Section~\ref{sec:impl-preferences}), and confirming, through twenty-run campaigns at
    $n=5{,}000$ and $n=10{,}000$ matching that same resolution, that this
    advantage is genuinely absent
    rather than merely unreproduced at larger scale: both confidence intervals
    are centred on zero and neither overlaps the $n=1{,}000$ result
    (Sections~\ref{sec:results-final-campaign} and~\ref{sec:results-scale}).

  \item An exploratory observation that the benefit of external sentence-level
    decomposition may be reduced when the underlying evaluator already performs
    comparable internal decomposition, offered as a hypothesis motivated by the
    mode comparison rather than as an established mechanism
    (Chapter~\ref{ch:discussion}, Section~\ref{sec:disc-interpretation}).

  \item An independent, held-out ground-truth evaluation (RQ4) showing that the
    $n=1{,}000$ advantage is not an artefact of validating against the trained
    models' own training-style labels: it holds, more sharply, against 500
    articles disjoint from every training set used in this thesis, and it
    reverses rather than merely vanishing as the training set grows, becoming a
    significant advantage for holistic scoring at $n=10{,}000$
    (Chapter~\ref{ch:results}, Section~\ref{sec:results-ground-truth}). Two
    candidate explanations for the reversal were tested directly and both
    rejected (Sections~\ref{sec:results-gt-margin}
    and~\ref{sec:results-gt-shortcut}).

  \item A transparent account of instability, stochastic variation across runs,
    methodological limitations, and the reproducibility of the pipeline
    (Chapters~\ref{ch:implementation}, \ref{ch:discussion},
    and~\ref{ch:conclusion}).
\end{enumerate}

These contributions concern the construction and learnability of preference data.
The thesis does not claim that GCA yields objectively more accurate factuality
labels, nor that it improves the factual consistency of generated summaries; the
scope of what the experiments support is stated precisely in
Chapter~\ref{ch:methodology}, Section~\ref{sec:methodology-scope}.

\section{Scope}
\label{sec:intro-scope}

\begin{itemize}
  \item \textbf{Task:} single-document abstractive news summarization, using the
    CNN/DailyMail dataset \citep{hermann2015teaching, nallapati2016abstractive}.
  \item \textbf{Candidate generation:} a single instruction-tuned base model,
    Mistral-7B-Instruct-v0.3, sampled at two temperatures ($T=0.7$ and $T=1.0$) to
    produce an A/B candidate pair per source article.
  \item \textbf{Judge:} a single fixed, deterministic factuality judge,
    \texttt{yzha/AlignScore} \citep{zha2023alignscore}, intended throughout
    in its \texttt{nli} evaluation mode, though the confirmatory campaigns
    were, due to an implementation oversight found during final review,
    actually run under its \texttt{nli\_sp} mode
    (Chapter~\ref{ch:implementation}, Section~\ref{sec:impl-preferences}).
    No generative LLM-as-judge and no external API dependency is
    used in the final pipeline.
  \item \textbf{Optimization target:} the learnability of the resulting preference
    data for Bradley--Terry reward-model training, measured by pairwise validation
    accuracy under 5-fold cross-validation. A downstream fine-tuned generation
    policy is not evaluated in the final comparison; an early Direct Preference
    Optimization (DPO) prototype was explored but dropped from the final scope
    (Chapter~\ref{ch:methodology}).
\end{itemize}

\section{Thesis Structure}
\label{sec:intro-structure}

Chapter~\ref{ch:related-work} situates this work within the literature on
preference-based alignment, LLM-as-a-judge reliability, fine-grained feedback for
long-form generation, and automatic factual-consistency evaluation.
Chapter~\ref{ch:background} introduces the technical mechanisms this thesis builds
on: the Bradley--Terry preference model, Direct Preference Optimization, the
AlignScore factuality metric, and the statistical tools
used to validate the results. Chapter~\ref{ch:methodology} presents the experimental
design, including the GCA aggregation formula itself, which is this thesis's own
contribution rather than prior work, and an explicit account of how the study that
was actually carried out diverges from the original proposal, and why. Chapter~\ref{ch:implementation}
describes the pipeline, codebase, and compute infrastructure. Chapter~\ref{ch:results}
reports the full experimental campaign, from the initial reward-model baseline
through the formula and judge-mode optimization to the final multi-seed and
multi-scale validation, and then to the independent ground-truth precision
evaluation that answers RQ4. Chapter~\ref{ch:discussion} interprets these
results, including why the ground-truth result reverses rather than merely
fading where the learnability result only fades, and
Chapter~\ref{ch:conclusion} answers the research questions, states the
limitations of this work, and outlines directions for future work.
